% =====================================================================
\section{Bounds for a global search}
\label{sm:prior}
% =====================================================================

This section supports Sec.~\ref{sec:global}: why the van Trees bound does not apply to a uniform window, the
Ziv--Zakai bound that replaces it, the two search bounds, the reach of a measurement and the best probe spacing,
the posterior at detection, and the conditions under which van Trees applies again after detection. Throughout,
$\bar\sigma_\theta^2(n)$ is the mean squared error of Eq.~\eqref{eq:bayes-mse}, averaged over the noise and over
$\theta$ uniform on the window of width $W$.

\paragraph{Why the window breaks van Trees.} For a prior density $p(\theta)$, the van Trees inequality
\cite{vantrees1968,gill1995} bounds the mean squared error of any estimator by
\begin{equation*}
  \bar\sigma_\theta^2(n)\ \ge\ \big(J(p)+nI_{\rm max}\big)^{-1},\qquad
  J(p)=\int\frac{p'(\theta)^2}{p(\theta)}\,d\theta .
\end{equation*}
A uniform window drops from $1/W$ to zero at its edges, so $J(p)$ is infinite and the inequality reduces to
$\bar\sigma_\theta\ge0$: true, but empty. The inverse variance $12/W^2$ cannot simply take its place, since
$J(p)\ge1/\mathrm{Var}(\theta)$ with equality only for a Gaussian. For a signal linear in $\theta$,
$(12/W^2+nI_{\rm max})^{-1/2}$ is the error of the best \emph{linear} estimator, which the best estimator beats.
Table~\ref{tab:sm-prior} gives the size of the effect, computed numerically for a measurement
$y=\theta+\varepsilon$ with $\varepsilon$ Gaussian of variance $1/(nI_{\rm max})$: the best estimator lies up to
$4.4\%$ below the additive curve, around $nI_{\rm max}W^2\approx100$, and the two agree at both ends.

\begin{table}[h]
\centering
\caption{Linear signal, uniform window. Ratio of two curves to the root of the best achievable mean squared error,
$\bar\sigma_{\rm best}$, as a function of the measured information $nI_{\rm max}W^2$. Additive: the van Trees form
with $J_0=12/W^2$, $(12/W^2+nI_{\rm max})^{-1/2}$. ZZ: the Ziv--Zakai bound, Eq.~\eqref{eq:sm-zz-linear}.}
\label{tab:sm-prior}
\begin{tabular}{lccccccc}
\hline
$nI_{\rm max}W^2$ & 1 & 10 & 30 & 100 & 300 & $10^3$ & $10^4$\\
\hline
additive$/\bar\sigma_{\rm best}$ & 0.999 & 1.009 & 1.032 & 1.044 & 1.036 & 1.024 & 1.009\\
ZZ$/\bar\sigma_{\rm best}$ & 0.932 & 0.923 & 0.956 & 0.980 & 0.990 & 0.995 & 0.998\\
\hline
\end{tabular}
\end{table}

\paragraph{The Ziv--Zakai bound.} The Ziv--Zakai bound \cite{zivzakai1969,bell1997} turns estimation into a set of
binary decisions. Take two positions a distance $h$ apart, $\theta$ and $\theta+h$, equally likely a priori, and let
$P_{\min}(\theta,\theta+h)$ be the smallest probability of a wrong decision that any test on the data can reach
between a dip at one and a dip at the other. Two steps give the bound. First, for the error $e=\hat\theta_n-\theta$
of any estimate,
\begin{equation*}
  \bar\sigma_\theta^2(n)=\mathbb E\big[e^2\big]=\frac12\int_0^\infty h\,\Pr\big(|e|\ge h/2\big)\,dh ,
\end{equation*}
which follows from $e^2=\int_0^{2|e|}(h/2)\,dh$, the average running over $\theta$ and the noise as in
Eq.~\eqref{eq:bayes-mse}. Second, the estimate is itself a test: decide for whichever of $\theta$ and $\theta+h$ lies
closer to $\hat\theta_n$. This test errs only if the estimate misses the true position by at least $h/2$, and no test
errs less often than $P_{\min}$, so
\begin{equation*}
  \Pr\big(|e|\ge h/2\big)\ \ge\ G(h)=\int\big[p(\theta)+p(\theta+h)\big]\,P_{\min}(\theta,\theta+h)\,d\theta .
\end{equation*}
Together,
\begin{equation}
  \bar\sigma_\theta^2(n)\ \ge\ \frac12\int_0^\infty h\,G(h)\,dh .
  \label{eq:sm-zz}
\end{equation}
Pairs that the data cannot tell apart thus add to the error in proportion to their distance. The original bound
also replaces $G$ by the smallest nonincreasing function above it, which can only raise the right-hand side; we do
not need it. The bound needs no derivative of $p$, so the edges of the window do not break it. Before any
measurement no test does better than a guess, $P_{\min}=\tfrac12$, and for the uniform prior $G(h)=(W-h)/W$, so
\begin{equation*}
  \bar\sigma_\theta^2(0)\ \ge\ \frac12\int_0^W h\,\frac{W-h}{W}\,dh=\frac{W^2}{12},
\end{equation*}
which is Eq.~\eqref{eq:prior}: the bound starts at the true prior error. For the linear signal,
$P_{\min}=Q\big(h\sqrt{nI_{\rm max}}/2\big)$ with $Q$ the Gaussian tail, which falls with $h$, and
\begin{equation}
  \bar\sigma_\theta^2(n)\ \ge\ \int_0^W h\Big(1-\frac hW\Big)\,Q\Big(\frac{h\sqrt{nI_{\rm max}}}{2}\Big)\,dh ,
  \label{eq:sm-zz-linear}
\end{equation}
which tends to $1/(nI_{\rm max})$ once $nI_{\rm max}W^2\gg1$, the data limit of Eq.~\eqref{eq:vantrees}. It lies within $8\%$ of the best achievable error at every $n$
(Table~\ref{tab:sm-prior}).

\paragraph{Search with a dip of finite reach.} Take the high-SNR idealization of the elimination model: a
measurement at $d$ reports the dip, without error, if and only if the dip lies within the probe cell of width
$\Delta d$ around $d$, and reports the baseline otherwise. Before its first detection every outcome of a run is the
baseline, whatever $\theta$ is, so any strategy, adaptive or not, has measured at a sequence $d_1,\dots,d_n$ fixed
in advance up to that point: each choice depends only on outcomes that are all the same. These $n$ probes cover a set
$C_n$ of measure at most $n\Delta d$; let $U$ be the rest of the window, of measure $L\ge W-n\Delta d$. If both
$\theta$ and $\theta+h$ lie in $U$, the data are the same under the two hypotheses and $P_{\min}=\tfrac12$. Keeping
only these pairs, with $p=1/W$,
\begin{equation*}
  G(h)\ \ge\ \frac2W\cdot\frac12\,\big|U\cap(U-h)\big|=\frac{\big|U\cap(U-h)\big|}{W},
\end{equation*}
where $|U\cap(U-h)|$ is the measure of the positions $\theta$ with both $\theta$ and $\theta+h$ in $U$, and
Eq.~\eqref{eq:sm-zz} gives
\begin{equation*}
  \bar\sigma_\theta^2(n)\ \ge\ \frac{1}{2W}\int_0^\infty h\,\big|U\cap(U-h)\big|\,dh
  =\frac{L^2}{4W}\,\mathbb E\,|X-Y| ,
\end{equation*}
with $X,Y$ independent and uniform on $U$. With $F$ their distribution function,
$\mathbb E|X-Y|=2\int F(1-F)\,dx$, and since $F$ grows by at most $dx/L$ over a step $dx$,
$\int F(1-F)\,dx\ge L\int_0^1u(1-u)\,du=L/6$, with equality only when $U$ is an interval. Hence
\begin{equation*}
  \bar\sigma_\theta^2(n)\ \ge\ \frac{L^3}{12\,W}\ \ge\ \frac{(W-n\Delta d)^3}{12\,W},
\end{equation*}
Eq.~\eqref{eq:search-mean}. The elimination model, which probes contiguously from one edge, meets
it: no strategy searches a window faster in this limit. A strategy that probes in a scattered order leaves $U$ in
pieces, and its mean error at the same $n$ is larger.

\paragraph{Search at any SNR.} At lower SNR a single measurement no longer decides, and the argument above no longer
applies, since an adaptive strategy now branches on noisy outcomes. A change-of-measure argument
\cite{garivier2019} bounds it instead. It rests on the same idea as the Ziv--Zakai bound, a binary test, but the
alternative to a dip at $\theta$ is now no dip at all, and the evidence for telling them apart is counted by the
Kullback--Leibler divergence. Let $P_\theta$ be the law of all the data of a run with the dip at $\theta$,
and $P_0$ the law with no dip at all, the strategy being the same. In each measurement the outcome is Gaussian with
variance $\sigma^2/k$ and mean shifted by the dip depth at the drive frequency,
$s(d;\theta)=C\gamma^2/[\gamma^2+(d-\theta)^2]$. Each drive frequency $d_i$ is chosen from the past outcomes by the
same rule under both laws, so the chain rule of the Kullback--Leibler divergence gives, for adaptive strategies too,
\begin{equation*}
  D_{\rm KL}(P_0\,\|\,P_\theta)=\mathbb E_0\Big[\sum_{i=1}^{n}\frac{k\,s(d_i;\theta)^2}{2\sigma^2}\Big].
\end{equation*}
Each term is the divergence of one measurement, for two Gaussians of equal variance the squared shift of the mean
over twice the variance. The expectation is under $P_0$, where the choices cannot depend on $\theta$. Averaged over
$\theta$ uniform on the window, at any $d$ one measurement adds at most
\begin{equation*}
  \begin{aligned}
  \bar\kappa&=\frac1W\int\frac{k\,s(d;\theta)^2}{2\sigma^2}\,d\theta
  \ \le\ \frac{k\,C^2\gamma^4}{2\sigma^2W}\int_{-\infty}^{\infty}\frac{dv}{(\gamma^2+v^2)^2}\\
  &=\frac{k\,C^2\gamma^4}{2\sigma^2W}\,\frac{\pi}{2\gamma^3}=\frac{\pi\gamma\,k\,\mathrm{SNR}^2}{4W},
  \end{aligned}
\end{equation*}
with $v=d-\theta$, the integral over the whole line bounding the one over the window. Hence
\begin{equation*}
  \frac1W\int D_{\rm KL}(P_0\,\|\,P_\theta)\,d\theta\ \le\ n\bar\kappa .
\end{equation*}
Now let $A_\theta$ be the event $|\hat\theta_n-\theta|\le t$, with probabilities $p_0(\theta)=P_0(A_\theta)$ and
$p_\theta=P_\theta(A_\theta)$. Under $P_0$ the estimate does not depend on $\theta$, so the average of $p_0(\theta)$
over the window is at most $2t/W$. By data processing and the joint convexity of the binary divergence
$\mathrm{kl}(p,q)=p\ln\frac pq+(1-p)\ln\frac{1-p}{1-q}$, the averages $\bar p_0$ and $\bar p$ satisfy
$\mathrm{kl}(\bar p_0,\bar p)\le n\bar\kappa$. With $\mathrm{kl}(p,q)\ge(1-p)\ln\frac{1}{1-q}-\ln2$, the fraction
of runs whose estimate is still more than $t$ from the dip is
\begin{equation*}
  1-\bar p\ \ge\ \exp\!\Big(-\frac{n\bar\kappa+\ln2}{1-2t/W}\Big),
  \qquad \bar\sigma_\theta^2(n)\ \ge\ t^2\,(1-\bar p).
\end{equation*}
This holds for every $0<t<W/2$. With $a=1+\ln2+n\bar\kappa$, the choice $t=W/2a$ gives
$1/(1-2t/W)=a/(a-1)$, so the exponent is $-a$ and
\begin{equation*}
  \bar\sigma_\theta^2(n)\ \ge\ \frac{W^2}{4a^2}\,e^{-a}=\frac{W^2}{8e}\,\frac{e^{-n\bar\kappa}}{(1+\ln2+n\bar\kappa)^2},
\end{equation*}
which is Eq.~\eqref{eq:search-bound}. It keeps the rate $e^{-n\bar\kappa}$ and holds at every $n$. Computed
numerically, it lies within a factor $1.03$ of the bound maximized over $t$ for $n\bar\kappa\le1$, and within $1.42$
for $n\bar\kappa\le50$. The fixed choice $t=W/4$ gives the simpler $W^2e^{-2n\bar\kappa}/64$, but at half the rate.
The search therefore takes at least of order $1/\bar\kappa$ measurements, Eq.~\eqref{eq:search-cost}. The bound has
no factor $\ln(W/\Gamma)$, and none is needed. A strategy that averages a fixed number of shots at every probe
position must raise its detection threshold with the number of positions to keep false alarms rare, and pays about
$\ln(W/\Delta d)$ more, Eq.~\eqref{eq:sm-fixed-cost} below. A sequential test spends its shots where they are
needed: it drops a position as soon as it looks like the baseline and measures again only where the outcome is
ambiguous, so it needs strong evidence only to confirm a dip, and avoids most of the factor \cite{malloy2014}. A
Bayesian strategy goes further: its posterior combines the outcomes at neighboring positions, which the tails of the
dip reach as well, and lowers the weight of a position with each outcome that looks like the baseline instead of
excluding it at once. At high SNR $1/\bar\kappa$ falls below $W/\Delta d$, one measurement per probe position, and the
finite-reach bound above is the one that binds.

\paragraph{The reach of a measurement.} The finite-reach bound treats the dip as invisible beyond its probe cell.
The probe spacing $\Delta d$ must make that true: a dip anywhere in a cell, even at its edge, a distance
$\Delta d/2$ from the probe, must be detected, while none of the $M=W/\Delta d$ cells on the baseline may be taken
for a dip. A probe declares a dip when its drop exceeds a threshold of $z_\alpha$ noise standard deviations
$\sigma/\sqrt k$. For Gaussian noise the largest of $M$ baseline draws is about $\sqrt{2\ln M}$, so keeping false
alarms rare needs $z_\alpha\approx\sqrt{2\ln M}$. A dip at the cell edge is
$C\gamma^2/[\gamma^2+(\Delta d/2)^2]=C/[1+(\Delta d/\Gamma)^2]$ deep, and is missed with a small probability only if
its drop exceeds the threshold by a further $z_\beta$. With $z=z_\alpha+z_\beta$,
\begin{equation}
  \frac{\sqrt k\,\mathrm{SNR}}{1+(\Delta d/\Gamma)^2}\ \ge\ z
  \quad\Longrightarrow\quad
  \Delta d=\Gamma\sqrt{\frac{\sqrt k\,\mathrm{SNR}}{z}-1}.
  \label{eq:sm-reach}
\end{equation}
A single measurement detects anything only for $\sqrt k\,\mathrm{SNR}>z$. The reach is below the line width for
$\sqrt k\,\mathrm{SNR}<2z$, where only the core of the dip stands out, and above it at higher SNR, where the
Lorentzian tails do too and a measurement rules out more than a line width. The change-of-measure bound uses the
Lorentzian itself and does not depend on this choice.

\paragraph{The best spacing.} At lower SNR a probe repeats $m$ measurements at the same drive frequency, which
multiplies $k$ by $m$ in Eq.~\eqref{eq:sm-reach}. More repeats widen the reach and so need fewer probes; the search
costs $m\,W/\Delta d$ measurements. With $a=\sqrt k\,\mathrm{SNR}/z$ and $u=\sqrt m\,a$, the cost per unit of window
is
\begin{equation*}
  \frac{m}{\Delta d}=\frac{1}{\Gamma\,a^2}\,\frac{u^2}{\sqrt{u-1}},
\end{equation*}
smallest at $u=4/3$, with $m$ treated as continuous. The best probe then averages just enough for the drop at
$\sqrt m\,\sqrt k\,\mathrm{SNR}=\tfrac43z$, and the best spacing is
\begin{equation*}
  \Delta d_{\rm opt}=\frac{\Gamma}{\sqrt3}\approx0.58\,\Gamma,
\end{equation*}
below the line width: it pays to see only the core of the dip and to take more probes. A spacing wider than
$\Gamma/\sqrt3$ is the best choice only when a single measurement already has $a>4/3$, so that $m=1$ and
$\Delta d$ follows from Eq.~\eqref{eq:sm-reach}. At the optimum the search costs
\begin{equation}
  n_{\rm fixed}=\frac{16\sqrt3}{9}\,\frac{W}{\Gamma}\,\frac{z^2}{k\,\mathrm{SNR}^2}
  \approx1.2\,z^2\,\frac{1}{\bar\kappa},
  \label{eq:sm-fixed-cost}
\end{equation}
with $1/\bar\kappa$ the change-of-measure bound, Eq.~\eqref{eq:search-cost}. With $z^2\approx2\ln M$ this is about
$2.4\ln(W/\Delta d)$ times the bound: the factor that a sequential test avoids.

\paragraph{The detection posterior.} Let the detecting measurement be at the center $d$ of its cell, with the dip
uniform over the cell; we neglect the weak information that the non-detection at the previous probe carries about the
near edge of the cell. The outcome $y$ of the measurement, and with it any detection rule applied to $y$, depends on
$\theta$ only through $|\theta-d|$. The posterior is then symmetric about $d$ and $\mathbb E[\theta\mid y]=d$ for
every $y$, so by the law of total variance
\begin{equation*}
  \mathbb E\big[\mathrm{Var}(\theta\mid y)\big]=\mathrm{Var}(\theta)-\mathrm{Var}\big(\mathbb E[\theta\mid y]\big)
  =\frac{\Delta d^2}{12}-0 ,
\end{equation*}
exactly, at any SNR. The depth of the drop changes only the shape of the posterior. At high SNR it fixes
$|\theta-d|$ to within about $I(\theta)^{-1/2}$, with $I(\theta)$ the Fisher information of Eq.~\eqref{eq:fisher}
of the detecting measurement, and the posterior is two narrow peaks at $d\pm|\theta-d|$.

\paragraph{Restarting van Trees after detection.} Let $\tau$ be the step at which a run declares the dip found,
decided from the run's own data and never from the true $\theta$. Given the data up to $\tau$, the posterior is the
prior of the remaining measurements, and van Trees applied to them, then averaged over runs, gives
$\bar\sigma_\theta^2(n)\ge\mathbb E[(J_\tau+(n-\tau)I_{\rm max})^{-1}]$ with $J_\tau$ the Fisher information of that
posterior. $J_\tau$ is finite: every likelihood is Gaussian in $y$ and so smooth in $\theta$, and the posterior has
none of the hard edges of the idealized cell. Right after detection the two mirror peaks make $J_\tau$ of order $I(\theta)$, far above
$12/\Delta d^2$, while the error is $\Delta d/\sqrt{12}$: the restarted bound is valid but loose until the side is
known. The elimination model therefore restarts with the value that matches the error, $J_\tau=12/\Delta d^2$, a
model value rather than a bound. The measurement that resolves the side is placed at the steepest slope of one of
the two peaks: if that is the true one it adds $I_{\rm max}$, as the model assumes; if not, it sees little of the dip
and adds less, but still tells the side. Once the side is known the posterior is a single peak of variance about
$(I(\theta)+I_{\rm max})^{-1}$, and the restarted bound is tight; the model, which keeps $12/\Delta d^2$ and
discards $I(\theta)$, is conservative from then on, by an amount that ranges from nothing for a dip right under the
probe, where $I(\theta)\approx0$, to $I_{\rm max}$ for one on its flank. The window edges still cut the posterior
off, but where it holds a negligible fraction of its mass, so smoothing them changes the bound by a negligible
amount. The bound is on the error expected given the run's history; the posterior variance of a single run may
fluctuate slightly below it.
